% ===== 原文 pp.10–12 =====
%
% tr-02 / chunk a —— 覆盖原文 p.10–p.12；公式 tag 2.31–2.49
% ---------------------------------------------------------------------

% [原文 p.10]
\subsection{Liouville 理论}

现在我们研究第二个二维场论，即 Liouville 理论。它是所谓非有理共形场理论的首要范例。
这里我们只涉及这一优美而丰富的理论中最浅显的方面。感兴趣的读者可参阅 [14,15] 等经典综述。
我们主要采用一套更偏公理化的进路，它由 [16] 开创并做了综述。在进入其中之前，
先让我们快速回顾一下自由玻色子理论。

\subsubsection{自由玻色子理论}

二维欧几里得号差（Euclidean signature）下的无质量玻色子 $\phi$ 满足运动方程
\begin{equation}
  \triangle\phi = \bar\partial\partial\phi = 0.
  \tag{2.31}
\end{equation}
于是通解为
\begin{equation}
  \phi(z,\bar z) = f(z) + \bar f(\bar z)
  \tag{2.32}
\end{equation}
其中 $f(z)$ 是一个全纯函数。

% [Tutor Note: p.10, 式 (2.31)：原文把二维拉普拉斯算子记作 $\triangle = \bar\partial\partial$，不写系数。
%   在 $z=x+\ii y$ 的约定下 $\triangle = \partial_x^2+\partial_y^2 = 4\partial_z\partial_{\bar z}$，故
%   $\triangle\phi=0 \Leftrightarrow \bar\partial\partial\phi=0 \Leftrightarrow \phi=f(z)+\bar f(\bar z)$，即 (2.32)。
%   原文省略的只是因子 4，不要以为 $\triangle$ 与 $\bar\partial\partial$ 之间还差一个待定的常数。]

由此我们看到，经典理论在坐标的任意全纯变换 $z\mapsto g(z)$ 下具有对称性。
把无穷小变换 $z\to z(1+\epsilon z^n)$ 的生成元记为 $L_n$。经典地，它们满足对易关系
$[L_m,L_n]=(m-n)L_{m+n}$。

在量子理论中，自由玻色子理论的基本算符乘积展开（OPE）是
\begin{equation}
  \partial\phi(z)\partial\phi(w) \sim -\frac{1}{2(z-w)^2}
  \tag{2.33}
\end{equation}
以及
\begin{equation}
  \ee^{2\ii a\phi(z,\bar z)}\ee^{-2\ii a\phi(w,\bar w)} \sim \frac{1}{|z-w|^{2a^2}}.
  \tag{2.34}
\end{equation}

生成元 $L_0$ 把坐标按 $z\mapsto \ee^{-\epsilon L_0}z$ 重新标度。相应地，当两点函数表现为
$(z-w)^{-2\Delta}(\bar z-\bar w)^{-2\bar\Delta}$ 时，我们就说该算符的全纯维数为 $L_0=\Delta$、
反全纯维数为 $\bar L_0=\bar\Delta$。我们发现 $\partial\phi$ 的 $(L_0,\bar L_0)=(1,0)$，
而 $\ee^{2\ii a\phi}$ 的 $(L_0,\bar L_0)=(a^2,a^2)$。

% [Tutor Note: p.10, "生成元 L_0 把坐标按 $z\mapsto \ee^{-\epsilon L_0}z$ 重新标度"：
%   原文把 $\epsilon L_0$ 直接放进指数，这里的 $L_0$ 要按它在该场上的本征值（即 $\Delta$）来理解；
%   通常的写法是坐标作 $z\mapsto \ee^{-\epsilon}z$ 的标度变换、其生成元记作 $L_0$。两种读法给出的
%   维数关系相同，但不要把 $\ee^{-\epsilon L_0}$ 当成作用在坐标上的算符去展开。]

% [Tutor Note: p.10, 式 (2.35)：中心项 $c(m^3-m)/12$ 在 $m=0,\pm1$ 时为零（$m^3-m=m(m-1)(m+1)$），
%   所以 $L_0,L_{\pm1}$ 生成的正是没有中心扩张的整体共形（SL(2,R)）子代数，中心荷只影响 $|m|\ge 2$ 的模。
%   换句话说，从经典对易关系 $[L_m,L_n]=(m-n)L_{m+n}$ 只能确定中心项的系数形状，$c$ 的取值必须靠
%   (2.36) 那样的 OPE 才能定出来（下页给出 $c=1$、$c=1+6Q^2$ 两个例子）。]

% [Tutor Note: p.10, 式 (2.33)--(2.34)：本页约定两点函数作 $(z-w)^{-2\Delta}(\bar z-\bar w)^{-2\bar\Delta}$，
%   且明说 $\ee^{2\ii a\phi}$ 的 $L_0=\Delta=a^2$，故它的两点函数应为 $|z-w|^{-4a^2}$。
%   由 (2.33) 的高斯收缩（$\langle\phi\phi\rangle=-\half\ln|z-w|^2$）也得到
%   $\langle \ee^{2\alpha\phi(z)}\ee^{2\beta\phi(w)}\rangle\propto|z-w|^{-4\alpha\beta}$：
%   取 $\alpha=\ii a,\beta=-\ii a$ 得 $|z-w|^{-4a^2}$；取 $\alpha=Q/2+\ii a,\beta=Q/2-\ii a$ 得 $|z-w|^{-4a^2-Q^2}$。
%   而原文 (2.34) 印作 $|z-w|^{-2a^2}$、(2.40) 印作 $|z-w|^{-2a^2-Q^2/2}$，恰好都只有一半，
%   疑为原文漏因子 $2$（已 650--900 dpi 裁图复核印刷原文，确实如此）。正文照原样排印；
%   读者若要与 (2.41)/(2.66) 的 $\Delta$ 对齐，请按 $|z-w|^{-4\Delta}$ 理解这两式的指数。]

在量子层面，由 $L_n$ 生成的代数被修正为
\begin{equation}
  [L_m,L_n] = (m-n)L_{m+n} + c\frac{m^3-m}{12}\delta_{m,-n}
  \tag{2.35}
\end{equation}
其中 $c$ 是一个称为中心荷的数。这就是著名的 Virasoro 代数。我们把这些生成元打包成一个场
$T(z)=\sum L_n z^{n-2}$。上面的对易关系等价于算符乘积展开
\begin{equation}
  T(z)T(w) = \frac{c}{2}\frac{1}{(z-w)^4} + \frac{2}{(z-w)^2}T(w) + \frac{1}{z-w}T'(w) + \cdots .
  \tag{2.36}
\end{equation}

% [原文 p.11]
在自由玻色子理论中，能量–动量张量 $T(z)$ 由 $T(z)=\half\partial\phi\partial\phi$ 给出。
稍作计算即可发现它满足关系 (2.36)，且 $c=1$。现在我们考虑一个略有修改的自由玻色子理论，
其能量–动量张量为
\begin{equation}
  T(z) = \partial\phi\partial\phi + \ii Q\partial^2\phi.
  \tag{2.37}
\end{equation}
我们发现它满足关系 (2.36)，其中
\begin{equation}
  c = 1 + 6Q^2.
  \tag{2.38}
\end{equation}

% [Tutor Note: p.11, 首句与式 (2.37)：首句写自由理论的 $T=\half\partial\phi\partial\phi$，
%   而 (2.37) 的第一项系数是 1（没有 $\half$）。两处不是同一归一化的直接比较：$c=1+6Q^2$ 是与 (2.37)
%   这一写法配套的结果，所以不要把 (2.37) 与首句的 $T$ 直接相减来"读"出 $Q$ 项的系数。]

这一修改对应于取拉格朗日密度
\begin{equation}
  \cL = g^{ij}\partial_i\phi\partial_j\phi + Q\phi R
  \tag{2.39}
\end{equation}
其中 $R$ 是二维曲面的曲率。在球面上计算关联函数时，我们可以把它映到平直的无限平面，
但要注意曲率仍然集中在 $z=\infty$ 处。这等效于在 $z=\infty$ 处放置一个算符 $\ee^{-2Q\phi}$。
出于这个原因，参数 $Q$ 常被称为背景电荷。这会把指数场的基本关联函数 (2.34) 修改为
\begin{equation}
  \ee^{2(Q/2+\ii a)\phi(z,\bar z)}\ee^{2(Q/2-\ii a)\phi(w,\bar w)}
  \sim \frac{1}{|z-w|^{2a^2+Q^2/2}}.
  \tag{2.40}
\end{equation}

特别地，指数算符 $\ee^{2(Q/2+\ii a)\phi}$ 具有 $L_0=a^2+Q^2/4$。更一般地，我们说算符
$\ee^{2a\phi}$ 具有
\begin{equation}
  L_0 = a(Q-a).
  \tag{2.41}
\end{equation}

% [Tutor Note: p.11, 式 (2.40)–(2.41)：把 $Q=0$ 代入 (2.40) 即回到 (2.34)，说明两式是同一约定下的
%   同一件事；再把 $a=Q/2+\ii p$ 代入 (2.41) 的 $a(Q-a)$，得 $(Q/2+\ii p)(Q/2-\ii p)=p^2+Q^2/4$，
%   正是 (2.40) 紧接一句里的 "$L_0=a^2+Q^2/4$"，也正是下页 (2.66) 里的 $\Delta_p$。
%   三处的 $L_0$ 可以互相对上；但注意 (2.40) 的 $|z-w|$ 指数本身比该约定少因子 $2$（见 (2.34) 处批注），
%   与 $\Delta_p$ 配套的是 $|z-w|^{-4\Delta_p}$——不要把 (2.40) 的指数直接当作 $2\Delta_p$ 使用。]

\subsubsection{相互作用理论}

现在假设我们想把这个带背景电荷 $Q$ 的自由玻色子理论变成一个完整的相互作用理论，
同时不破坏共形不变性。在领头阶，拉格朗日量里的新项应当具有维数 $(L_0,\bar L_0)=(1,1)$，
才能保持共形不变性。由 (2.41) 我们发现，只要
\begin{equation}
  Q = b + \frac{1}{b},
  \tag{2.42}
\end{equation}
算符 $\ee^{2b\phi}$ 恰好能完成这件事。现在拉格朗日密度为
\begin{equation}
  \cL = g^{ij}\partial_i\phi\partial_j\phi + Q\phi R + 4\pi\mu\ee^{2b\phi}
  \tag{2.43}
\end{equation}
其中参数 $\mu$ 常被称为宇宙学常数。

% [Tutor Note: p.11, 式 (2.42)：这个条件的全部用处就是让 $\ee^{2b\phi}$ 的共形维数恰为 1——由 (2.41)
%   取 $a=b$ 得 $L_0=b(Q-b)=b\cdot(1/b)=1$（同理 $\bar L_0=1$）。于是 $4\pi\mu\ee^{2b\phi}$
%   是一个 $(1,1)$ 型的边际形变，可以加进拉格朗日量而不破坏共形不变性。注意反过来读：
%   (2.42) 是把"要一个 $(1,1)$ 算符"这个要求解出来的，不是额外假设。]

这个参数 $\mu$ 可以通过平移 $\phi$ 的原点而被取为任意想要的值，因此以这个相互作用项为基础
做微扰论并不容易。换个说法，势项 $\ee^{2b\phi}$ 是指数型的，不能看成对具有拉格朗日量 (2.39)
的自由理论的一个小形变。但经过一系列令人印象深刻的工作，我们现在知道
% [原文 p.12]
拉格朗日密度 (2.43) 决定了一个相互作用的共形场理论，其中心荷为
\begin{equation}
  c = 1 + 6Q^2 = 1 + 6\left(b+\frac{1}{b}\right)^2.
  \tag{2.44}
\end{equation}

% [Tutor Note: p.12, 式 (2.44)：$c=1+6(b+1/b)^2$ 在 $b\to b^{-1}$ 下不变，这正是下页讨论的
%   "$b$ 反演不变性"在中心荷层面的体现；再结合 $Q=b+1/b$ 也在此变换下不变，就可以理解为什么
%   两点、三点函数在 $b\to b^{-1}$ 下本质不变。该不变性在 (2.43) 的拉格朗日描述里完全看不到
%   （经典作用量里 $b$ 与 $b^{-1}$ 地位并不对称），这是 Liouville 理论的一个量子效应。]

与自由理论相比，有一个关键差别，如下所述。在自由理论中，按态–算符对应，具有 $L_0=p^2$ 的
算符 $\ee^{2\ii p\phi}$ 给出一个沿 $\phi$ 所参数化的方向以动量 $p$ 运动的态。特别地，
$\ee^{2\ii p\phi}$ 与 $\ee^{-2\ii p\phi}$ 这两个算符是两个不同的态。

类似地，在 Liouville 理论中，具有 $L_0=p^2+(Q/2)^2$ 的算符 $\ee^{(Q+2\ii p)\phi}$ 对应于
$\phi$-空间中动量为 $p$ 的态。这里把指数平移 $Q$ 是必要的，为的是让 $L_0$ 保持实且正。
由于指数型相互作用 $\ee^{2b\phi}$，来自 $\phi$ 负值区域的波无法穿透到 $\phi$ 正值区域，
而是被指数势垒反射。这意味着 $\ee^{(Q+2\ii p)\phi}$ 与 $\ee^{(Q-2\ii p)\phi}$ 这两个算符
给出同一个态：
\begin{equation}
  \ee^{(Q+2\ii p)\phi} = R(p)\ee^{(Q-2\ii p)\phi}
  \tag{2.45}
\end{equation}
其中 $R(p)$ 是一个称为反射系数的相位。

% [Tutor Note: p.12, 式 (2.45) 上面那句"平移 $Q$ 是必要的，为的是让 $L_0$ 保持实且正"：
%   直接算一下就看出来——由 (2.41)，$\ee^{2a\phi}$ 的 $L_0=a(Q-a)$。若用不平移的 $\ee^{2\ii p\phi}$
%   即 $a=\ii p$，得 $L_0=\ii p(Q-\ii p)=p^2+\ii pQ$，只要 $Q\neq0$ 就是复数；而取
%   $a=Q/2+\ii p$ 即 $L_0=p^2+Q^2/4\ge Q^2/4>0$，实且正。这正是"必须平移 $Q$"的全部算术内容。]

$R(p)$ 的渐近行为可以用该指数势对应的量子力学散射问题算出；而共形不变性对 $R(p)$ 的
高阶修正所施加的约束，正是完整求解 Liouville 理论的出发点之一。

我们现在知道，一个酉共形理论可以由如下条件唯一确定：它的谱由一族主场 $V_p(z)$ 给出，
其中 $p$ 为实数且 $p\ge 0$，都具有 $L_0=\bar L_0=p^2+(Q/2)^2$。中心荷是 (2.44)。
球面上的三点函数最初由 [17] 与 [18] 各自独立得到。利用共形不变性，我们可以把一个算符放在
$z=\infty$、另一个放在 $z=1$、第三个放在 $z=0$。于是它由下式给出
\begin{equation}
  \begin{aligned}
    C(\alpha_1,\alpha_2,\alpha_3) &:= \langle \ee^{2(Q-\alpha_1)\phi(\infty)}
      \ee^{2\alpha_2\phi(1)}\ee^{2\alpha_3\phi(0)}\rangle \\
    &= \left[\pi\mu\gamma(b^2)b^{2-2b^2}\right]^{(Q-\alpha_1-\alpha_2-\alpha_3)/b} \\
    &\quad \times \frac{\Upsilon'(0)\Upsilon(2\alpha_1)\Upsilon(2\alpha_2)\Upsilon(2\alpha_3)}
      {\Upsilon(\alpha_1+\alpha_2+\alpha_3-Q)\Upsilon(\alpha_1+\alpha_2-\alpha_3)
       \Upsilon(\alpha_1-\alpha_2+\alpha_3)\Upsilon(-\alpha_1+\alpha_2+\alpha_3)}.
  \end{aligned}
  \tag{2.46}
\end{equation}

% [Tutor Note: p.12, 式 (2.46)：注意 $z=\infty$ 处的算符写成 $\ee^{2(Q-\alpha_1)\phi}$ 而不是
%   $\ee^{2\alpha_1\phi}$。原因是 $h_\alpha=\alpha(Q-\alpha)$（见 (2.66)）在
%   $\alpha\leftrightarrow Q-\alpha$ 下不变，$\ee^{2\alpha\phi}$ 与 $\ee^{2(Q-\alpha)\phi}$
%   的共形维数相同；把 $Q-\alpha$ 放在 $\infty$、把 $\alpha$ 放在 $0$ 只是约定。
%   (2.57) 里的 $C(\alpha_1,\alpha_2,\alpha)$ 与 (2.55) 里的 $\langle\cdots\ee^{2\alpha_4\phi}(0)\rangle$
%   都沿用同一约定，抄写时不要"顺手"把它们改成 $2\alpha_1\phi$。]

这里，
\begin{equation}
  \gamma(x) = \Gamma(x)/\Gamma(1-x)
  \tag{2.47}
\end{equation}
以及
\begin{equation}
  \Upsilon(x) = \frac{1}{\Gamma_2(x|b,b^{-1})\Gamma_2(Q-x|b,b^{-1})}
  \tag{2.48}
\end{equation}
其中 $\Gamma_2(x|\epsilon_1,\epsilon_2)$ 是 Barnes 双 Gamma 函数，它由正规化如下无穷乘积得到
\begin{equation}
  \Gamma_2(x|\epsilon_1,\epsilon_2) \propto \prod_{m,n\ge 0}(x+m\epsilon_1+n\epsilon_2)^{-1}.
  \tag{2.49}
\end{equation}

% [Tutor Note: p.12, 式 (2.48)–(2.49)：由 (2.48) 立得 $\Upsilon(x)=\Upsilon(Q-x)$（分母两个因子互换），
%   这是后面所有 DOZZ 表达式的基本对称性。再配合 (2.50) 的复共轭性质：$b$ 实、$x$ 实（此时
%   $Q=b+1/b$ 也实）时 $\Upsilon(x)$ 为实数，故 (2.46) 的四个 $\Upsilon$ 因子在实参数下都是实的
%   （前面那个带幂次的因子另计）。这是用 DOZZ 公式做数值时最基本的自检。
%   另注意 (2.49) 只定义到相差一个与 $x$ 无关的整函数因子，原文用 $\propto$ 正是为此；
%   后面所有公式里 $\Upsilon'(0)$ 之类的整体常数都与这个约定绑定。]
